\documentclass[letterpaper]{article}
\usepackage[submission]{aaai2027}
% Fonts are managed by the unmodified official AAAI-27 style.
\usepackage[hyphens]{url}
\usepackage{graphicx}
\urlstyle{rm}
\def\UrlFont{\rm}
\usepackage{natbib}
\usepackage{caption}
\frenchspacing
\usepackage{amsmath,amssymb,booktabs,array}
\pdfinfo{/TemplateVersion (2027.1)}
\setcounter{secnumdepth}{0}
\newtheorem{proposition}{Proposition}
\newenvironment{proof}{\par\noindent\textit{Proof.}\ }{\hfill$\square$\par}
\newcommand{\enc}{\Phi}

\title{Compositional Narrative Planning through Source-Preserving Domain Fusion\\Supplementary Material}
\author{Anonymous Authors}
\affiliations{}
\begin{document}
\maketitle
\section{Scope and Reading Guide}
This supplement gives the optional method details, the learned-domain first-repeat evaluation, its expanded 70-task diagnostic, the separate 11-step authored case, and dated earlier development evidence. Their different sources, budgets, and evaluation units are not pooled into a success rate. The first repeat is an interrupted continuation; composition tasks are separately constructed worlds. Source execution, interface semantics, and narrative quality remain distinct claims. Working review questions at the end expose unresolved evidence.

\section{Additional Method Details}
The main paper describes how domains are connected and how source actions are retained. This section records an optional reduction for a fixed world and the generic search implementation used in earlier development studies. The learned-domain evaluation uses the frozen groups and Glaive settings reported in the main paper.

\subsection{Optional Reduction for a Fixed World}
For a supplied world, let $U_i$ contain the objects satisfying source $i$'s immutable root-membership predicate. Every declared source constant must retain its recorded memberships, and every positive subtype membership must imply root membership. This admission check matters because a constant literal in an effect has no action-parameter guard. The selector rejects missing constant memberships rather than inferring or restoring facts.

Construct a graph over sources. Sources sharing a target predicate are adjacent if that predicate is nullary or their root universes overlap. A positive-arity interface cannot communicate a tuple between disjoint root universes, regardless of its argument permutation. Let $S$ be the union of graph components containing every declared owner of a positive or negative goal atom. The reduction retains all actions from $S$, together with every original declaration, constant, and the exact mapped problem. It does not recompile a smaller object universe. Its certificate records the source universes and retained owners; reuse with changed universes or an unsupported goal owner requires recertification.

\begin{proposition}[Fixed-universe reduction]
Assume the stated membership compilation, immutable memberships covering source constants, source-owned action variants without additional adapters, and a ground goal whose atoms lie in their declared source universes. Let $D_S$ remove only the actions outside $S$. Then the supplied goal is reachable in $D_S$ if and only if it is reachable in the full compiled domain.
\end{proposition}
\begin{proof}
Membership-guarded parameters and quantifiers, together with constant coverage, restrict each source's effective predicate reads and writes to tuples from its own $U_i$. A nullary shared predicate creates an edge. For a positive-arity shared predicate, an excluded source and a retained source have disjoint universes, so their possible tuples are disjoint. Excluded actions therefore leave all retained source-tuple truth values unchanged. The truth of retained conditions and their conditional effects depends only on those tuples and immutable memberships; declarations and quantifier universes remain identical. Deleting excluded actions from a full solution consequently preserves each retained action's applicability and the supplied goal. Conversely, every reduced solution is a full-domain solution because retained actions and the problem are unchanged.
\end{proof}

\paragraph{Boundary.}
The statement concerns goal reachability, not full-state trace equivalence: excluded independent actions can change other tuples or events. It provides no minimum-component, bounded-search success, narrative-quality, or interface-semantic guarantee, and it does not apply to initial-only outcome ranking. The proof argument and its relationship to existing modularity and decomposition work still require independent technical review. A deliberately malformed mapped-world control exposes a lost producer when constant coverage is omitted; the strengthened admission check rejects that input. Rechecking all thirty valid component inputs in the matched source-world block leaves their generated domains byte-identical.

\subsection{Generic Search and Outcome Generation}
\paragraph{Goal-directed mode.}
The default reference search uses bounded BFS and static-positive parameter projection. The latter restricts candidate domains only with conjunctive static Boolean atoms having false defaults; remaining correlations and dynamic conditions are checked by the simulator. State, depth, grounding, and subprocess time limits are reported distinctly. The optional Pyperplan reference uses greedy best-first search with hFF after recorded UP compilations; its adapter enforces only a subprocess time budget. Every returned plan is independently checked in the original exported model after compiler plan lifting.

\paragraph{Initial-only mode.}
Without a goal, bounded BFS returns the first distinct noninitial states in discovery order. Each candidate records its positive and negative changes from $I$, a nonempty plan, and validation against the corresponding generated endpoint condition. This is a declared outcome-generation policy, not narrative ranking or success against independently specified goals. Resource cutoffs remain visible even when some outcomes have been generated.

For false-default Boolean states, the implementation maintains the complete set of true ground atoms. Conjunctive positive preconditions project parameter choices against this current-state index; remaining conditions are checked by simulation. All expanded effect targets update the index from the successor state, including deletions and conditional or quantified effects. Endpoint differences are the symmetric difference of the initial and successor true-atom sets. Grounding and fluent budgets count unique visited projected bindings and effect targets on visited state-changing transitions, respectively. These counters differ from eager enumeration. Authored finite controls agree with unpruned Cartesian BFS on result order, plans, and complete changes; they do not establish corpus-wide scaling.


\section{Learned-Domain First-Repeat Evaluation}
\label{sec:supp-repeat-one-evaluation}

We evaluate whether a merged narrative planning domain preserves executable source plans and supports new plans that connect actions from different source domains. These are distinct requirements: retaining an original plan does not establish cross-source composition, and finding a composite plan does not establish preservation of every source behavior. We therefore separate a frozen source-preservation experiment from a subsequent, explicitly constructed composition diagnostic. Compression and original-goal search coverage are reported as descriptive properties, not as requirements for successful composition.

\subsection{Data, conditions, and execution protocol}

\paragraph{Source domains and sampling.}
We use the GPT-4 revised domain--problem artifacts from \emph{There and Back Again}~\cite{kelly2023there}, with the upstream version pinned in the audit record below. The recorded cohort contains 192 published records: 40 previously designated development records and 152 legacy test records. We retain the previously admitted 60-source parser/compiler cohort, comprising 32 TinyStories and 28 WritingPrompts domains; the remaining 92 legacy test records remain documented exclusions. No source-domain repair or new admission decision is performed for this experiment. Because these sources have already been used during development, the legacy split does not establish unseen-source generalization.

The frozen dataset contains 100 balanced source sequences of length ten. We evaluate their prefixes of size $k\in\{2,3,4\}$, yielding 100 groups per size and 300 groups in total. At each size, every source appears either $\lfloor100k/60\rfloor$ or $\lceil100k/60\rceil$ times, and unordered groups are unique within a size. Sampling uses source identifiers rather than semantic labels, alignment decisions, or planning outcomes. The resulting groups contain 538 distinct source pairs. Groups share sources and are nested within sequences, so they are correlated evaluation units rather than 300 independent stories.

\paragraph{Conditions.}
\emph{Native} searches each original domain--problem pair. \emph{Union} combines source domains while retaining separate predicate and action namespaces. \emph{Predicate-only} adds accepted predicate correspondences while retaining source action variants. \emph{Full} uses exactly the same predicate correspondences and additionally merges eligible action variants with identical mapped effects and acting roles, using disjoined, source-guarded preconditions. Full is evaluated as a domain-composition method; it is not required to outperform Predicate-only on reachability. All 900 parameterized merged domains and their provenance mappings are saved before main-experiment planning.

\paragraph{Alignment and planning.}
The frozen candidate inventory contains at most eight predicate and four action candidates per source pair, totaling 3,970 predicate and 710 action judgments. Candidates are deduplicated within the repeat. GLM-5.3 Cloud is queried through Ollama with temperature 0, requested seed 2026110101, a 32,768-token context setting, and a 16,384-token generation setting. Provider enforcement of the requested seed is unverified. Correspondences require confidence at least 0.9, valid role permutations and source-local evidence. Multi-source clusters additionally require complete pairwise compatibility, consistent role cycles, and agreement on concept, scope, and aligned role descriptions. Invalid or unavailable judgments do not authorize sharing.

Each search and replay has a 15-second limit. Glaive runs with one planner worker and a 4-GB Java memory limit. We use the separately versioned runtime that corrects false-disjunction substitution and grounding over empty types; historical results from the uncorrected runtime are excluded. The main experiment contains 60 Native searches and 1,800 merged-domain searches: two fixed original source problems per group, under each of three merged-domain conditions. Every source in every group also has a registered original-plan transport slot, totaling 2,700 slots.

\paragraph{Repeat status.}
This report analyzes the first completed repeat, retained as an \emph{interrupted continuation}. Its serial and ten-client request segments, failed responses, and one administratively interrupted request are preserved; unsuccessful requests are not replaced. Five repeats were registered, but only the first is complete. The second stopped after 576 attempts when the service reported its monthly usage limit, and is excluded from the completed-repeat metrics. We consequently report no across-repeat standard deviation, reproducibility estimate, or significance test. We also avoid confidence intervals that would incorrectly treat overlapping source/group rows as independent observations.

\subsection{Evaluation units and validity criteria}

\paragraph{Source-plan preservation.}
Let $\mathcal{S}$ be the registered source-plan/domain slots for a condition. Let $E_i$ indicate that slot $i$ has a Native plan passing both Glaive replay and an independent physical state-transition check. Let $V_i$ indicate that its translated plan passes the same checks in the merged domain. We report
\[
 P_{\mathrm{available}}=\frac{\sum_{i\in\mathcal{S}}E_iV_i}{\sum_{i\in\mathcal{S}}E_i},
 \qquad
 C_{\mathrm{evidence}}=\frac{\sum_{i\in\mathcal{S}}E_i}{|\mathcal{S}|}.
\]
Missing Native plans remain unavailable, rather than being counted as successful transports or preservation failures. Preservation is evaluated on explicitly mapped original initial states and goals. Each problem activates only its designated source; facts from unrelated stories are not pooled. This tests individual source-world embeddings, not the consistency of arbitrary joint worlds or preservation of all possible traces.

\paragraph{Returned plans.}
A returned plan is validated only when Glaive reports goal achievement with zero unexplained steps and the independent physical checker confirms applicability and goal achievement. We distinguish validated plans, unsupported/unknown returns, bounded searches with no returned plan, and time limits. A bounded no-plan result is not a proof of unsolvability. Checking an output schema or replaying a plan does not provide an independent semantic label.

\paragraph{Cross-source composition.}
For the separate diagnostic, a plan counts as a witnessed composition only if its goal-relevant causal slice contains an edge from an action of one source to an action of another source through a shared predicate. Re-adding an initially true fact does not create such an edge; unrelated actions concatenated before a goal do not count. Ambiguous provenance of a fused action is classified as unknown rather than automatically credited to multiple sources. This measures witnessed composition in the returned plan, not the existence or absence of every possible composite plan.

\subsection{Alignment outcomes and merged artifacts}

The model produced contract-valid responses for 4,117 of 4,680 candidates (87.97\%), of which 206 proposed compatibility (Table~\ref{tab:judgments}). These are automatic judgments, not accuracy measurements. The 561 invalid outputs comprise 389 source-citation mismatches, 96 JSON parsing failures, 69 incomplete/truncated generations, and seven incomplete compact judgments. Two further outcomes are unavailable: a transport failure and an administratively interrupted in-flight request. All remain in the candidate denominator.

\begin{table*}[htbp]
\centering\small
\caption{Automatic judgment outcomes. Contract-valid responses are not semantic correctness labels; compatible is a subset of contract-valid.}
\label{tab:judgments}
\begin{tabular}{lrrrrr}
\toprule
Kind & Candidates & Contract-valid & Compatible & Invalid & Unavailable \\
\midrule
Predicate & 3,970 & 3,426 & 182 & 543 & 1 \\
Action & 710 & 691 & 24 & 18 & 1 \\
Total & 4,680 & 4,117 & 206 & 561 & 2 \\
\bottomrule
\end{tabular}
\end{table*}


Predicate sharing occurs in 120 of 300 groups, and action merging in 19 groups (Table~\ref{tab:models}). Across correlated group occurrences, the compiler accepts 169 predicate clusters and 19 action clusters, representing 95 and 11 distinct member sets, respectively. All accepted clusters in this repeat are pairwise. Eleven of the 19 action-cluster occurrences have different business preconditions while retaining identical mapped effects and acting roles. Thus, this run contains actual action-merge cases, although their occurrence is sparse and is not itself an endpoint of the composition evaluation.

\begin{table*}[htbp]
\centering\small
\caption{Structural outcomes by group size. Predicate and action totals count schema occurrences across groups; reductions are descriptive, not primary success criteria.}
\label{tab:models}
\begin{tabular}{lrrrrr}
\toprule
$k$ & Groups & With sharing & With action merge & Predicates U$\to$F & Actions U$\to$F \\
\midrule
2 & 100 & 15 & 0 & $1,074\to 1,057$ & $877\to 877$ \\
3 & 100 & 42 & 7 & $1,610\to 1,556$ & $1,315\to 1,308$ \\
4 & 100 & 63 & 12 & $2,151\to 2,053$ & $1,751\to 1,739$ \\
Total & 300 & 120 & 19 & $4,835\to 4,666$ & $3,943\to 3,924$ \\
\bottomrule
\end{tabular}
\par\smallskip\begin{minipage}{0.97\linewidth}\small With sharing counts groups with accepted predicate correspondences. U: Union; F: Full. The 300 groups reuse 60 sources.\end{minipage}
\end{table*}


The action-merge cases do not provide exhaustive branch coverage. Of 38 registered source-branch occurrences in the 19 action clusters, 19 have an available validated Native plan, and 11 are actually exercised by a fused step in a transported Native plan. We rechecked the effect/actor equality conditions for all 19 cluster occurrences; these structural checks do not supply independent evidence of semantic equivalence.

\subsection{Preservation and source-local search results}

All available source plans survive translation into Full. Native search returns 34 plans out of 60 problems: 33 validate, while one contains unsupported counterfactual steps and remains unknown. Sixteen Native searches return no plan within their execution and ten hit the wall-time limit. The 33 validated Native plans supply 499 eligible Full transport slots across the overlapping groups; all 499 validate (Table~\ref{tab:preservation}). The remaining 401 of 900 Full slots are unavailable, giving 55.44\% evidence coverage rather than universal empirical coverage.

\begin{table*}[htbp]
\centering\small
\caption{Original-plan transport in Full. Evidence coverage includes unavailable slots in its denominator; all three merged conditions have identical counts.}
\label{tab:preservation}
\begin{tabular}{lrrrr}
\toprule
$k$ & Source slots & Retained / available & Unavailable & Evidence coverage \\
\midrule
2 & 200 & 111/111 & 89 & 55.50\% \\
3 & 300 & 165/165 & 135 & 55.00\% \\
4 & 400 & 223/223 & 177 & 55.75\% \\
Total & 900 & 499/499 & 401 & 55.44\% \\
\bottomrule
\end{tabular}
\par\smallskip\begin{minipage}{0.97\linewidth}\small 499 available slots reuse 33 distinct validated Native plans. No preservation violation was observed, but unavailable plans are not validated.\end{minipage}
\end{table*}


Union and Predicate-only have the same transport outcomes as Full. Across all three conditions, 1,497 available slots validate and 1,203 are unavailable; these totals repeatedly reuse the same 33 Native plans and must not be interpreted as independent plan samples. Within the 120 groups with actual predicate sharing, Full validates 206 available transports out of 408 registered slots. Within the 19 groups with actual action merging, it validates 35 available transports out of 69 slots. Neither subset contains an observed transport violation, but unavailable plans and unexercised branches remain evidence gaps.

Fresh source-local search outcomes are identical across the three merged-domain conditions on all 600 matched anchor queries. Per condition, 341 of 600 searches return a plan; 332 validate and nine remain unknown because of counterfactual plan syntax. There are 160 bounded no-plan outcomes and 99 wall-time limits (Table~\ref{tab:search}). Across conditions, the 27 unknown returns all concern repeated uses of \texttt{tiny\_stories\_\_story\_094}. These results are regression evidence: by construction, only one source is active per query. They neither test cross-source generation nor establish that composition is ineffective.

\begin{table*}[htbp]
\centering\small
\caption{Fresh source-local searches, per merged condition. All 600 matched query statuses are identical across Union, Predicate-only, and Full.}
\label{tab:search}
\begin{tabular}{lrrrrrr}
\toprule
$k$ & Queries & Returned & Validated & Unknown & No plan & Time limit \\
\midrule
2 & 200 & 114 & 111 & 3 & 54 & 32 \\
3 & 200 & 114 & 111 & 3 & 53 & 33 \\
4 & 200 & 113 & 110 & 3 & 53 & 34 \\
Total & 600 & 341 & 332 & 9 & 160 & 99 \\
\bottomrule
\end{tabular}
\par\smallskip\begin{minipage}{0.97\linewidth}\small Each query activates one original source. No plan denotes a bounded result, not certified unsolvability. Nine unknown returns per condition have unsupported counterfactual steps.\end{minipage}
\end{table*}


\subsection{Historical 18-task composition diagnostic}

We evaluated composition in a separately registered posthoc diagnostic using the already constructed first-repeat domains. A restricted generator considers positive producer--consumer predicate interfaces with conjunctive supported preconditions and without constant bindings. It specifies objects, cross-source identity/type assignments, initial facts, intentions, and goals explicitly; it copies no source-problem facts. The original parameterized merged domains remain byte-identical. The generated worlds are structural mechanism fixtures, not original story worlds or independent test stories.

The generator produces a supported task for 35 of the 300 groups: six, twelve, and seventeen at sizes two, three, and four. The other 265 groups are classified as having no task generated by this restricted rule, not as composition failures. Before search, we select the first six generated groups per size under a fixed hash ordering. This gives 18 tasks with 13 distinct explicit worlds. Selection is conditional on accepted interfaces and generator support, but never on search success. We run Union, Predicate-only, Full, and three Full controls that disable the producer source, disable the consumer source, or separate the shared object identities, for 108 searches in total.

Full returns 18 validated plans, of which 17 exhibit a goal-relevant cross-source causal link (Table~\ref{tab:composition}). Predicate-only has the same result, establishing that shared predicates with retained action variants suffice for these witnessed compositions. The remaining Full plan reaches its goal through a single source; this does not prove that a composite alternative is absent. The 17 composite plans span 12 distinct worlds, all have two steps, and all actually use two sources. The counts by selected domain size are 6/6, 6/6, and 5/6; they do not establish three- or four-source causal chains.

\begin{table*}[htbp]
\centering\small
\caption{Posthoc composition diagnostic: 18 selected tasks comprising 13 distinct explicit worlds. A composite plan requires a goal-relevant cross-source causal link.}
\label{tab:composition}
\begin{tabular}{lrrr}
\toprule
Condition & Validated plans & Composite plans & No plan \\
\midrule
Union & 3/18 & 0/18 & 15 \\
Predicate-only & 18/18 & 17/18 & 0 \\
Full & 18/18 & 17/18 & 0 \\
Full, producer disabled & 3/18 & 0/18 & 15 \\
Full, consumer disabled & 0/18 & 0/18 & 18 \\
Full, separate identities & 3/18 & 0/18 & 15 \\
\bottomrule
\end{tabular}
\par\smallskip\begin{minipage}{0.97\linewidth}\small The 17 Full composite plans represent 12 distinct worlds and all use two sources. These are conditional diagnostic counts, not success rates on the 300-group sample.\end{minipage}
\end{table*}


Fifteen selected tasks have a validated composite Full plan while Union and all three controls return no plan within the fixed budget. This is bounded control evidence for the tested interface/identity assumptions, not a proof that the goal is unreachable without either source. In two further tasks, Full returns a cross-source plan even though a single-source alternative also reaches the goal. These remain witnessed compositions: the objective is to create valid combinations, not to require an increase in original-goal solution coverage.

\paragraph{Example.}
For \texttt{story\_028} and \texttt{story\_085}, an explicit fixture declares a common character $r_0$ and item $r_1$, initial fact $\mathrm{reads}(r_0)$, and the character's intention and goal $\mathrm{playing}(r_0,r_1)$. Under the closed-world convention, the character initially lacks the item and the item is not broken. The planner returns
\[
\mathrm{find}_{028}(r_0,r_1)
\;\to\;\mathrm{has}_{\mathrm{shared}}(r_0,r_1)
\;\to\;\mathrm{play}_{085}(r_0,r_1).
\]
The arrows show the causal handoff, not a third action. Both physical replay and Glaive's explanation/goal checks pass. Turning off predicate sharing, disabling either source, or separating the identities produces no returned plan in the registered controls. Natural-language realization can narrate these two events, but narrative quality was not evaluated.

\paragraph{Preservation in the same diagnostic models.}
For the 18 selected Full models, we additionally register all 54 source slots and translate their available Native plans. All 27 available transports validate; the other 27 slots remain unavailable. Fifteen of the 17 composite groups have at least one preserved Native plan, and two groups have available, preserved plans for every selected source. These two groups share one synthetic world. That world identifies a lightbox with the object occupying the \texttt{?bike} role of a pedal-repair action; its broad formal type permits the binding, but its semantic plausibility is unverified. We retain this case as a concrete identity concern, not an independently established semantic error or a clean generalization success.

\subsection{Interpretation and limits}

The completed repeat supports two bounded findings: available source traces are preserved in their original-world embeddings, and the fixed merged domains can produce formally valid cross-source plans under explicitly declared composition worlds. The experiment does not establish preservation of all possible source traces, compatibility of arbitrary pooled worlds, correctness of automatic object identities, or narrative coherence. Output validation and Glaive explanation checks are not substitutes for independent semantic assessment.

The sample and execution also limit statistical interpretation. The source cohort was admitted by an existing parser/compiler contract; groups overlap and nest, and diagnostic tasks repeat explicit worlds. Only one interrupted repeat is complete, with mixed request concurrency and a changing service execution history. We therefore make no five-repeat stability, unseen-source, timing-speedup, or population-level significance claim. The historical 17/18 count and the expanded counts below are conditional diagnostics, not estimates over all 300 groups.

The next evidence needed for the stated objective is broader source-trace availability, independently checked cross-source identity mappings, and fixed tasks requiring longer causal chains involving three or four sources. Additional repetitions can quantify alignment variability after service capacity is restored. None of these gaps is repaired by treating unavailable rows as successes, expanding the same two-event template, or reporting structural compression as the main outcome.

\subsection{First-Repeat Audit and Reproduction}
The read-only recomputation decoded all 4,678 available raw responses and reproduced their recorded output-contract results. Two unavailable outcomes remain in the 4,680-candidate denominator. It rechecked 525 unique source-symbol contexts, consensus for all 300 groups, all 900 compiled models, and the effect/actor certificates of all 19 action-cluster occurrences.

Physical rechecks cover 2,526 previously validated instances: 33 Native plans, 996 returned merged-domain plans, and 1,497 transports. They contain 1,245 distinct domain/problem/plan triples and 9,132 step checks, with no detected discrepancy. These are correlated rechecks, not new independent experiments. The 27 unsupported merged returns and one unsupported Native return remain unknown. All 2,700 original-world problems activate only the designated source. No new model request or planner search was made for the audit or manuscript integration.

The upstream source commit is
\begin{quote}\ttfamily
de333cbc07a2c7359425\\
16bf6ef87d55d4b723a2
\end{quote}
The frozen dataset, continuation protocol, corrected Glaive binary, and model digest are respectively identified by the following SHA-256 values; displayed line breaks are not part of the hashes.
\begin{quote}\ttfamily
80bf350c90e68e44ddfeb2e4c807b4bc\\
4ee5ac4f4c117cbdf44c11342df544df
\end{quote}
\begin{quote}\ttfamily
5f8f0a677a0c817b89a6ed1750f07961\\
5fe9630066c945d732abef23a2475bfe
\end{quote}
\begin{quote}\ttfamily
0e2df21e54471634bb5087882f311e281\\
8d9b9764eec726addbbb8ce9b18e225
\end{quote}
\begin{quote}\ttfamily
8477dab3e25bb0f93c468af220186f55\\
394262ee5e9f39262af4b60b54a8c4ba
\end{quote}

The first-repeat audit retains group/source metrics, branch certificates, raw-response receipts, physical-recheck records, and a file manifest. The 3 November ICAPS evidence inventory verifies 233 pinned reports: 227 historical records and six first-repeat assessment reports. This heterogeneous inventory count is neither a statistical sample size nor an aggregate success denominator. It does not replace the earlier dated snapshots. The local source package contains the manuscript sources and an evidence-location manifest; the full experiment remains in the repository pending an anonymized release.

\section{Expanded Composition Diagnostic and Long Case}
\label{sec:supp-expanded-composition}

\paragraph{Coverage and units.}
The first-repeat 300-group sample is unchanged. Of its 120 groups with accepted predicate sharing, 83 contain no positive producer-effect/consumer-precondition pair in the generator's supported action fragment, two have such a pair but no distinct positive consumer goal effect, and 35 support task construction. The other 180 groups have no accepted predicate interface. Thus 35/300 measures the coverage of this particular construction rule, not a merge success rate. Seventy-four source actions are excluded for non-STRIPS formulas and three for constant bindings; these action-level exclusions are not an exclusive decomposition of the 85 uncovered predicate-sharing groups. Missing interfaces may reflect genuine incompatibility, candidate-retrieval limits, conservative judgments, or invalid responses; independent semantic labels are absent.

The expanded study retains every supported descriptor: 8, 26, and 36 tasks from 6, 12, and 17 groups at input sizes two, three, and four. Seventy tasks span 37 exact world keys, formed from objects, facts, goals, and intentions. Different merged-domain contexts may share a world key. For the primary deduplicated report, one representative is fixed before search by smallest input size and then seeded case hash. The remaining contexts are retained; no best-of-outcome selection is used. Distinct world keys still share source models and are not independent samples.

\paragraph{Fresh execution.}
The new matrix contains 420 searches and 212 returned-plan validations. Each call retains the original 15-second, 4-GB, single-worker budget. Inputs, model copies, mappings, identities, goals, and implementation hashes are frozen before search; no new LLM requests or retries occur. All 212 returned plans pass original-domain Glaive explanation/goal checks and independent physical replay, and their physical and causal checks are recomputed offline. The old 18 descriptors are included within the 70; all 108 overlapping search, validity, and causal statuses reproduce. They are not extra independent tasks or a second LLM merge repeat.

\begin{table*}[t]
\centering
\caption{All expanded outcomes. C: confirmed goal-relevant cross-source link; U: valid plan with origin-ambiguous causal classification; N: valid plan with no such link; P: no plan returned within budget. Each block retains its full task denominator.}
\label{tab:expanded-deduplicated}
\begin{tabular}{lrrrrrrrrrr}
\toprule
& \multicolumn{5}{c}{All 70 descriptors} & \multicolumn{5}{c}{37 fixed world representatives}\\
\cmidrule(lr){2-6}\cmidrule(lr){7-11}
Condition & Valid & C & U & N & P & Valid & C & U & N & P\\
\midrule
Union & 25 & 0 & 0 & 25 & 45 & 13 & 0 & 0 & 13 & 24\\
Predicate-only & 68 & 56 & 0 & 12 & 2 & 36 & 30 & 0 & 6 & 1\\
Full & 68 & 50 & 10 & 8 & 2 & 36 & 27 & 5 & 4 & 1\\
Full, producer off & 25 & 0 & 10 & 15 & 45 & 13 & 0 & 5 & 8 & 24\\
Full, consumer off & 1 & 0 & 0 & 1 & 69 & 1 & 0 & 0 & 1 & 36\\
Full, identities split & 25 & 0 & 10 & 15 & 45 & 13 & 0 & 5 & 8 & 24\\
\bottomrule
\end{tabular}
\end{table*}

\paragraph{Origin ambiguity.}
All ten Full unknowns have the reason \texttt{ambiguous\_fused\_source\_attribution}. Six are confirmed cross-source under Predicate-only and four have no cross-source link there. The other 60 classifications agree. The same ambiguity affects ten producer-off and ten identity-split results; zero confirmed links does not establish that all connections disappear. Six of the 35 models actually merge actions, covering 17 descriptors. In that subset, both Predicate-only and Full validate 16 plans; the former confirms 11 combinations, while Full confirms five and leaves ten unknown. This is a provenance limitation, not an established loss of executable behavior. All 50 confirmed Full plans use two actions and two sources.

\paragraph{Original-plan availability.}
The same 35 saved Full models have 116 source--group slots. Historical, hash-linked checks validate all 59 available original plans; 57 slots remain unavailable. No new original-plan calls are made. Of 50 confirmed composition descriptors, 47 have at least one preserved original plan; five have an available preserved reference for every selected source, spanning three worlds. The existing lightbox/bike concern occurs in three descriptors sharing one world and remains disclosed.

Separately, the 33 sources with validated Native plans induce 57 fully reference-covered groups within the original frozen sample: 31/16/10 at sizes 2/3/4. All 150 original-plan checks pass; 12 groups share predicates and two merge actions. Only six expanded composition descriptors belong to fully covered groups. This is a posthoc availability analysis, not a replacement sampling protocol or a completed new 33-source experiment. The 33 Native plans have lengths 1--20; 18 exceed two steps. Original-plan replay therefore includes multi-step traces even though the automatic composition fixtures remain two-step.

\subsection{An Authored Eleven-Step Case with Four Sources}
The existing Full model for group \texttt{k04\_0218ded14bc9db22} retains all 19 source actions: A=\texttt{story\_028}, B=\texttt{story\_085}, C=\texttt{story\_063}, and D=\texttt{wp\_060}. No source rule is repaired. We explicitly declare Mira, Noah, Emma, a toy, a walnut, and a coat with source memberships. Mira belongs to A/B/D, Noah to B/C, Emma to C, the toy to A/B, the walnut to C, and the coat to D. Initial ordinary facts are A.\texttt{reads(mira)}, C.\texttt{hard(walnut)}, D.\texttt{alive(mira)}, and D.\texttt{has(mira,coat)}; other facts are false under CWA. No original problem facts are pooled.

The conjunctive goal records Mira as cleaned, Noah as having forgiven Mira, the walnut as eaten, Emma as helping Noah, and Mira as wearing the coat. Initial intentions explicitly cover these goals. The final authored version also states Mira's intention to apologize. Table~\ref{tab:four-source-long-plan} preserves the actual planner order.

\begin{table}[t]
\centering
\caption{The separately authored four-source plan. Source labels identify the unchanged original actions.}
\label{tab:four-source-long-plan}
\begin{tabular}{rcl}
\toprule
Step & Source & Ground action\\
\midrule
1 & C & help(emma,noah)\\
2 & B & apologize(mira,noah)\\
3 & D & adjust-item(mira,coat)\\
4 & D & wear-item(mira,coat)\\
5 & B & forgive(noah,mira)\\
6 & C & break-object(emma,walnut,noah)\\
7 & C & eat-together(emma,noah,walnut)\\
8 & A & find(mira,toy)\\
9 & B & play(mira,toy)\\
10 & B & break-item(mira,toy)\\
11 & B & clean(mira,toy)\\
\bottomrule
\end{tabular}
\end{table}

Step 8 establishes shared possession for step 9, followed by playing--breaking--cleaning. The other chains are 2$\to$5, 1$\to$6$\to$7 (also 1$\to$7), and 3$\to$4. All four sources and all steps are in the goal-relevant causal slice. Forty-four action pairs have neither a causal path nor read/write interference; independence means separate goal branches, not unexplained filler. This is four-source participation through one cross-source edge plus local branches, not a single chain traversing all four sources.

\paragraph{Search and revision history.}
The original authored world receives seven direct Glaive searches: 60 seconds for Full and 15 seconds for Union, four source-off controls, and separated identities. All seven time out; they establish no control necessity result. A separate Fast Downward 26.6 LAMA-first search finds 12 physical steps, but UP rejects an empty optional \texttt{:types} declaration in the search export. A separate export omits that empty declaration and uses \texttt{seq-opt-lmcut}, returning the 11-step plan with independent physical validation. For physical search only, intentional annotations are omitted; all action parameters, preconditions, effects, physical facts, and goals are unchanged. Validation returns to the original intentional model.

Both physical candidates initially replay to the goal but have one unexplained/orphan step, the apology. A diagnostic identifies that step; in a separately frozen revision we add only Mira's explicit intention to apologize. The original merged domain, physical initial facts, goal, and 11 actions remain byte-identical. Original-domain Glaive replay then has zero unexplained and orphan steps, with independent physical replay also passing. The old worlds, timeouts, parser error, and failed intention checks are retained. This is an authored development revision, not a blind benchmark success or another model-judgment repeat. Source B still lacks a validated original-problem reference plan, so the case does not establish preservation of all four original plans. Narrative quality and source-model semantics are not independently assessed.

\section{31 October Development Evaluation}
We organize the development evidence around three questions: whether interfaces enable a causal chain requiring both sources, whether available source executions survive composition, and what action compression adds when predicate mappings are fixed. The illustrative archive example in the main paper's conceptual figure is not part of these experiments. The measured composition case in this historical block is the acquire-then-handoff family below; narrative interest and broad plot diversity are not measured.

\subsection{Inputs, Conditions, and Reproducibility}
This historical development pool comprises 15 known narrative domain/problem pairs with 234 actions and 236 predicates. Upstream learner versions and training provenance are not established for this pool. We retain original worlds and all failed or limited tasks. Explicitly authored composition worlds, original-goal problems, and schema witnesses are different evaluation units. Repeated representations, plans, and sizes of one authored family are correlated observations.

The 31 October 2026 block freezes source bytes, candidate lists, prompts, thresholds, model settings, and planner slots before execution. Candidate retrieval uses typed, same-named effect predicates and complete parameter bijections. It is a restricted opportunity census, not a semantic-recall evaluation. The primary confidence threshold is 0.9. Conditions are no shared predicates, predicate sharing with source variants, and the same predicate sharing plus accepted action merges. Source-removal controls retain the supplied world and vocabulary while removing the designated source's actions.

Model queries use \texttt{glm-5.3:cloud} on 31 October 2026, temperature 0, seed 20261027, context limit 32,768 and output limit 16,384 tokens, one attempt and one worker, with a 240-second request limit and a 900-second overall budget. Prompts provide source formulas and role/event evidence, request separate predicate/action judgments with citations, and exclude task goals, plans, and labels. The manifest records the model digest; a cloud identifier is not a guarantee of immutable server weights. The supplementary record specifies the digest and frozen prompt hashes. Invalid responses are retained without retries.

Planning uses Fast Downward 26.6 with LAMA-first~\cite{helmert2006downward}, a 20-second per-call wall budget, 4096 MiB memory, and one worker. Unified Planning 1.3.0 independently validates returned plans. Direct forward transports receive 30 seconds each and check applicability and goal satisfaction; they do not enumerate all source-atom truth values. The experiments ran on a shared Windows workstation. A complete run-linked hardware and operating-system receipt remains a reproducibility gap; concurrent workloads also preclude isolated speed claims. Artifact hashes identify the frozen implementations and exported models.

\subsection{Candidate Funnel and Action Compression}
All 105 pairs of the 15 sources remain in the census. Of 234 actions, 225 enter the candidate fragment and nine are outside it. The coarse scan finds 626 pairs across 80 source pairs, narrowed to seven typed action candidates. Nine predicate candidates and seven action candidates produce 16 model requests: 12 outputs pass the contract checker and four are invalid. The action decisions comprise three compatible, two different, and two invalid outputs. Only the previously known BIRT/SOUP travel pair reaches the fixed 0.9 threshold; the gift candidates score 0.85 and 0.88 and remain separate.

Six unique predicate pairs are accepted, represented as seven groups across eight full conditions because two gift cases reuse a predicate judgment. These are not independent semantic labels. Travel compresses its two-source model from 24 actions to 23; no newly discovered action-merge positive is established. The final block contains 25 compiled models and the complete 96-call planning denominator shown in Table~\ref{tab:story-summary}.

\begin{table}[t]
\centering
\begin{tabular}{lr}
\toprule
Outcome or evaluation unit & Count\\
\midrule
Planner: solved & 25\\
Planner: reported unsolvable & 58\\
Planner: wall-time limit & 13\\
\midrule
Original-plan transport: valid & 19\\
Original-plan transport: unavailable & 44\\
Separate schema witnesses: valid & 42\\
\bottomrule
\end{tabular}
\caption{Complete denominators for the 31 October development block: 96 planner calls and 105 transport slots. The 19 original-plan transports cover four unique original problems; witnesses are separate authored units. Unavailable plans are not successful preservation checks.}
\label{tab:story-summary}
\end{table}

\subsection{Preservation and Composition Controls}
Native planning supplies four original-goal plans from 15 problems; eight are reported unsolvable and three reach the time limit. Across the fixed representations, all 19 available original-plan transports and all 42 separate schema witnesses validate, totaling 216 transported steps without a detected violation. All 25 newly returned plans also revalidate. The 44 slots without a native plan remain unavailable. Three checks specifically exercise the merged action: BIRT's original 12-step trace and a one-step witness for each travel branch. SOUP's original two-step solution does not invoke travel and alone would not test that branch.

An authored acquire-then-handoff family connects PUSS and REMA actions. At one, four, and sixteen items, both predicate-only and full conditions solve in 2, 8, and 32 steps. No-share and either-source-removal controls are reported unsolvable at all three sizes (Table~\ref{tab:handoff}). These are three sizes of one development family, not three independent domains. The source \texttt{buy\_boots} action permits creating possession without payment or item-specific restrictions. We retain this underspecification; mechanical source necessity does not demonstrate realistic purchasing or narrative quality.

\begin{table}[t]
\centering
\begin{tabular}{lrrr}
\toprule
Condition & 1 item & 4 items & 16 items\\
\midrule
No share & U & U & U\\
Predicate-only & 2 & 8 & 32\\
Full & 2 & 8 & 32\\
Remove PUSS & U & U & U\\
Remove REMA & U & U & U\\
\bottomrule
\end{tabular}
\caption{Returned plan lengths in one authored composition family. U denotes solver-reported unsolvability. Predicate sharing accounts for the observed reachability gain; full composition adds no success in this family.}
\label{tab:handoff}
\end{table}

\subsection{Historical Controls and Evidence Boundaries}
Earlier frozen preservation bundles address different questions and budgets. A 60-second, 4-GiB LAMA-first run supplies six original-goal plans among 15 sources; all six translate into the same fixed union and zero-merge membership models. A separate action-invocation bundle retains all 234 target schemas: 146 valid prefixes, 50 reported unsolvable targets, and 38 wall-time limits. All 146 returned prefixes replay in both fixed models, with 354,757 mapped source-atom comparisons per representation and no mismatch. Invocation probes replace the goal with an observer flag; they are not original-goal successes or complete conditional-effect coverage. The supplementary evidence ledger preserves these distinct denominators and dated conditions.

A matched 75-call comparison on the same 15 original worlds finds four native successes. Full-domain zero-merge and fixed-interface conditions each solve zero, while static-component restriction recovers four in each condition. Thus that block establishes no source-goal benefit from its fixed interfaces. Separately, fresh reference constructions end at their three-hour and five-hour limits without a final root. A cached condition is a different experiment and fails final validation. Unavailable reference models are not counted as planner failures, and none of these incomplete constructions provides a completed efficacy baseline.

Authored finite-state fixtures exercise permutations, type universes, quantified formulas, conditional effects, and variable capture. Mutation controls show that the checker detects deliberately introduced errors. A separate semantic-negative example equates being captured alive with being dead: both isolated source traces remain valid, yet the false correspondence enables a capture--bury plan. Such controls establish mechanism behavior and checker sensitivity; they cannot estimate real-world semantic accuracy. Full details and retained failure rows appear in the supplement.


\section{31 October Model and Artifact Record}
The model identifier is \texttt{glm-5.3:cloud}; the recorded digest is

\begin{quote}\ttfamily
8477dab3e25bb0f93c468af220186f55\\
394262ee5e9f39262af4b60b54a8c4ba
\end{quote}.
This digest identifies the recorded manifest, not immutable cloud weights.
The frozen packet SHA-256 is

\begin{quote}\ttfamily
d10f4f7c30d7e99e975f7f90af270578\\
fe300b5e6236e8023d97afd6d6385bac
\end{quote};
the final requests SHA-256 is

\begin{quote}\ttfamily
843953e833a50abdaae8a2a5aaf5be82\\
23e715f328d7bc5f185f92497e6806e1
\end{quote}.
The source artifacts retain all 16 requests and raw responses, 96 planner calls, 105 transport slots, and invalid or unavailable rows. An evidence-inventory rebuild on 2 November verifies 227 heterogeneous records; this is an inventory count, not a statistical sample size. Full prompts and machine-readable outcomes remain in the local artifact package and require anonymized release preparation.

\section{Historical Development Evidence through 24 October}
\paragraph{Data boundary.}
This historical pool contains 15 local narrative domain/problem pairs, 234 actions, and 236 predicates. Complete schema/problem recovery passes for the corrected source-preserving union. Upstream learner versions and training provenance have not been supplied for this pool. These models and previously examined examples are development material; they do not establish independent-domain generalization.

\begin{table*}[t]
\centering
\begin{tabular}{lrrrr}
\toprule
Completed source-search run & Valid plans & Incomplete & Limits & Errors\\
\midrule
Legacy BFS, 60 s/source & 0 & 0 & 10 & 3\\
Static-projection BFS, 60 s/source & 0 & 0 & 13 & 0\\
Pyperplan GBFS+hFF, 120 s/source & 4 & 8 & 3 & 0\\
Fast Downward LAMA-first, 60 s/source & 6 & 0 & 1 & 0\\
\bottomrule
\end{tabular}
\caption{Development coverage across 15 sources per row. The remaining two cases in each BFS row and eight Fast Downward cases were reported unsolvable. Legacy limits comprise four state limits and six time limits. Static-projection limits comprise five state limits, three grounding limits, and five timeouts; Pyperplan limits are timeouts. Fast Downward 26.6 uses LAMA-first with a 4 GiB memory budget; its one cutoff is a memory limit. Different engines and budgets prevent treating this table as a controlled method comparison.}
\end{table*}

\paragraph{Conditional preservation evidence.}
The original Pyperplan run supplies four source plans, all valid after translation. A subsequent Fast Downward 26.6 LAMA-first run~\cite{helmert2006downward} supplies six: ACCO (6 steps), BIRT (12), CHIC\_NP (40), HANS\_NP (8), MUSI (14), and SOUP (2). All six validate in the same fixed namespaced union and after zero-merge type-membership compilation. The current measured result is consequently $6/6$ among available source plans, with source availability $6/15$. The remaining sources comprise eight solver-reported unsolvable tasks and one memory cutoff (SNOW), not nine failed plan translations. Domain hashes remain unchanged during replay; no source-goal repair is performed.

\paragraph{Auxiliary source-effect probes.}
To broaden trace coverage without repairing the original goals, a separate frozen bundle contains three sampled positive effect targets per source (45 probes, seed 20261015). Full source domains, objects, and initial facts remain unchanged; only each probe goal is replaced. Fast Downward returns 26 valid nonempty plans and reports 19 probes unsolvable. All 26 available plans validate after translation into each of the same fixed union and zero-merge membership models, covering 13 source domains. A separate stepwise audit compares all mapped ground source atoms along each available trace: both models pass 52,116 truth comparisons across 73 states, without mismatches or atom-budget cutoffs. This covers these trajectories, not all reachable states. These auxiliary targets are development probes, not original-goal successes, independent domains, or semantic annotations. The complete denominator remains 45, and sources without an available probe trace remain uncovered.

\paragraph{Trace-coverage boundary.}
Fresh validation of the six original-goal plans and 26 auxiliary plans reveals 61 distinct source action schemas across 234 declared actions. Original plans cover 31 schemas and auxiliary plans cover 38, with overlap; 173 remain unvisited. The only source action declaring a conditional effect is unvisited. Visiting an action schema would still not establish coverage of all its bindings or conditional branches. Thus domain coverage and large ground-atom comparison counts must not be interpreted as comprehensive action coverage; quantified-effect implementation evidence currently comes from authored controls.

Separately authored HANS prefixes exercise the unvisited conditional action with its guard false and true, starting from the unchanged original initial state. Both prefixes execute, and 2,605 mapped-atom comparisons over five states agree after compilation. They are not original-goal solutions and do not change the frozen probe denominator or its reported coverage. The true-guard prefix uses an initially possessed pebble in an action named for dropping breadcrumbs: the source action accepts any item. This retained source-model generality limits narrative interpretation even when transition preservation succeeds.

\paragraph{Action-invocation probes.}
A further bundle freezes one task for each of the 234 source action schemas before search. For target action $a$, introduce a fresh nullary Boolean $m_a$, initially false, and conjoin $m_a$ to the action's original effect; no precondition or original effect subtree changes. Replace only the problem goal by $m_a$. Projection onto the original fluents preserves the initial state, action applicability, and every transition: the observer occurs in no original condition and writes no original fluent. Consequently, a solution supplies an executable original-world prefix containing an invocation of $a$, while any such prefix satisfies the monitor goal. This statement concerns action invocation, including when an original conditional-effect guard is false; it does not establish effect-branch coverage or attainment of the original goal. Each task receives a 20-second, 4-GiB LAMA-first budget. Returned traces require independent plan validation, target-occurrence checking, and full source-atom comparison before fixed-model replay. The complete run returns 146 independently valid plans, reports 50 tasks unsolvable, and reaches 38 wall-time limits. Instrumentation inversion and target-occurrence checks pass. All 146 prefixes execute in each fixed union and membership model; each comparison covers 459 states and 354,757 source-atom truth values with zero mismatches. No returned trace is omitted from these audits.

Across the invocation traces and the previously validated original-goal and auxiliary-goal traces, 166 distinct source actions have execution witnesses. Twenty of these witnesses resolve actions whose own invocation task timed out; the 146/50/38 target-task outcome counts remain unchanged. Fifty actions are solver-reported unreachable and 18 remain unresolved, with no witness contradicting a reported unsolvability result. Among the witnessed schemas, 67 declare delete effects and 15 declare negative conditions. The conditional-effect schema still has no witness in these three frozen solver bundles; the two separately authored branch controls above remain separate evidence. These counts characterize the supplied source worlds, not all possible instances of the domains.

\paragraph{Source-problem diagnostics.}
A conservative predicate-symbol relaxation provides independently inspectable reachability obstructions for three of the eight solver-reported unsolvable source problems (BUFF, JACK\_NP, PUSS). The diagnostic ignores argument bindings, negative conditions, universal preconditions, and conditional-effect guards, thereby enlarging the set of potentially reachable predicates. Failure to support the relaxed goal implies an obstruction; support does not establish concrete solvability. For example, BUFF requires \texttt{cloak\_shaken} in its goal, but no initial fact or action effect establishes that predicate. The other five unsolvability results are not explained by this coarse diagnostic. No initial facts or actions are added as repairs, and an unsolvable source problem does not imply that its parameterized domain is unusable with another explicitly supplied world.

\paragraph{External development validation.}
An admission inspection considers six original-format families from the Sabre benchmark archive~\cite{ware2026collection}. Basketball's ten problems require an explicit packaging adapter: add the missing domain-reference header and remove problem-level requirements already declared by the unchanged domain. All other problem sections and domain bytes are preserved. The remaining five inspected families are not admitted because their original files require additional declaration handling or unsupported intentionality semantics; no intention features are stripped. Fast Downward finds independently valid plans for all ten packaged Basketball tasks. Each plan also validates in a fixed 16-source namespaced union and a fixed zero-merge membership model. Each representation passes 38,560 mapped-atom comparisons over 80 states, with unchanged model hashes. These are ten problems from one external development family, not ten independent domains or a blinded semantic-generalization test. Each problem retains its own initial world; other sources' narrative facts are not imported.

\paragraph{Engineering controls.}
The saved test checkpoint contains 124 passing component tests. Additional persistent synthetic mode checks cover initial+goal, initial-only, goal-only rejection, an explicitly empty state, a second object, and explicit source selection. Synthetic alignment controls test permutation direction, type assumptions, shared-entity declarations, quantified universe preservation, and retention of all actions after domain selection. These controls are not an independent narrative benchmark.

A raw-declaration guard additionally rejects duplicate sections, symbols and variables before dictionary-based parsing can overwrite them. An authored duplicate-action input previously lost one action while its generated output passed parsing; it is now rejected before domain export. Rebuilding the 15 local sources after this input-validation change yields byte-identical compiled domain and problem PDDL to the preceding fixed compiler version. Legal variable shadowing in separate scopes remains supported.

\paragraph{Matched-world comparison.}
An additional frozen comparison uses intact ACCO and BIRT action models, three explicitly authored worlds, two interface conditions, and three selection policies (18 conditions total). Both interface conditions use identical source-membership encoding and shared entities; they differ only in whether an authored possession correspondence is present. Pair-only and dependency selection retain the same two domains here, so these are not independent retrieval successes. All-domain selection retains all 15 domains. Every search receives a 60-second Pyperplan subprocess budget. Six plans validate independently, six searches are incomplete, and six reach the time limit. Input, model and source hashes pass the post-run audit.

\begin{table*}[t]
\centering
\begin{tabular}{llll}
\toprule
Authored world & Zero merge: pair/dep. & Interface: pair/dep. & All: zero/interface\\
\midrule
Initial BIRT location & I / I & S4 / S4 & T / T\\
Initial ACCO possession & S3 / S3 & S3 / S3 & T / T\\
Empty narrative state & I / I & I / I & T / T\\
\bottomrule
\end{tabular}
\caption{All 18 matched development conditions. S$k$: independently valid $k$-step plan; I: incomplete search; T: time limit. No I or T is classified as proven unsolvable. The worlds and interface are authored hypotheses, not independently annotated test cases.}
\end{table*}

\paragraph{Same-export stronger-planner check.}
Fast Downward 26.6 with LAMA-first subsequently solves the same 18 exported PDDL tasks under a 60-second wall budget and an additional 4 GiB memory limit. It returns nine valid plans and reports the other nine tasks unsolvable; every returned plan passes fresh independent validation. The audit checks byte-for-byte equality of both domain and problem inputs across solvers. All three authored-interface bridge conditions solve in four steps; all six source-ready controls solve in three steps. The three zero-merge bridge conditions and six empty-world conditions are reported unsolvable. Thus the Pyperplan cutoffs do not establish failure of the full-domain representation. This is a solver-coverage result on authored development cases, not a narrative-quality comparison or a controlled runtime speedup.

\paragraph{Interpretation limit.}
The four-step bridge plan uses one BIRT action to obtain possession, then three ACCO actions to arrange a performance and gain wealth. The BIRT model allows the first action to produce possession of an arbitrary typed item; this source-model generality is retained and disclosed, not repaired using the desired outcome. The example therefore demonstrates execution through an explicit interface, not narrative plausibility. Concurrent workstation workloads also prevent treating these durations as controlled performance measurements. A stronger planner comparison must reuse the exact exported domain/problem bytes.

\paragraph{Source-action necessity control.}
For each of the six explicit-pair world/interface conditions, a further frozen comparison retains all actions, removes all ACCO actions, or removes all BIRT actions (18 tasks). Vocabulary, entity memberships, initial facts, goals, and interfaces remain unchanged. Fast Downward returns five independently valid plans and reports 13 tasks unsolvable, without resource cutoffs. The authored-interface bridge is solved with both sources but reported unsolvable after either action removal. Both source-ready conditions remain solvable after BIRT action removal and become unsolvable after ACCO action removal. All empty-world and zero-interface bridge conditions remain unsolvable. Independent structural comparison confirms that only the declared actions changed. This establishes action-source necessity for the specific authored bridge task under the recorded solver result, not a general property of the domain pool or evidence of narrative quality.

\paragraph{Initial-only comparison.}
Removing the goal from the six explicit-pair world/interface conditions yields a separate complete initial-only run. All six searches reach the 60-second wall limit with no returned outcomes, while frozen source, query, and implementation hashes remain unchanged. Sparse effect-target comparison alone therefore does not establish a practical outcome-generation method on these cases. A bounded stack-sampling diagnostic locates ongoing work in eager action grounding; this is an engineering diagnosis, not a controlled performance attribution.

A subsequent current-state projection run uses byte-identical queries and PDDL, with the same 60-second wall budget and separately documented visited-binding counters. Each of the four nonempty-world conditions returns three outcomes; the two empty-world conditions exhaust their reachable state space without an outcome. Independent revalidation confirms all 12 plans and their complete endpoint differences. An omitted-change negative control remains plan-valid but fails the difference audit. These results support coverage of the declared outcome-generation policy on the six development cases, not narrative quality or a controlled runtime speedup.

\paragraph{Intermediate reference diagnostic.}
An earlier fresh adapter run terminated with a worker error after 10,285.9 seconds during a host interruption, with 6,460 cached responses and 13 completed fusion nodes. Source snapshots and implementation hashes remained unchanged; there is no completed final root. A separately frozen completed ACCO/BIRT fusion node maps 29 source actions to 28 actions. Under the adapter's positional problem export and unchanged ordered action formals, the available ACCO source trace remains valid, whereas the BIRT trace fails at step five: the rewritten action adds \texttt{has(nora,key)} as a precondition, which the original action instead establishes as an effect. The first four steps execute in both models. This is a trace-preservation counterexample for that intermediate node, not the final reference result, evidence of unsolvability, or a judgment that the original action is more narratively plausible.

A later frozen intermediate frontier consists of three disjoint completed nodes covering all 15 sources once. Of 234 source actions, 166 match after positional predicate/constant renaming and 68 do not. All three target domains parse independently. Each of the six available source plans encounters an action without verified parameter transport: across 14 distinct encountered mappings, ten change arity, one changes ordered types, and three change formal names or order. The last category may include benign renaming; it is not classified as execution failure. This frontier is still incomplete fusion, and these mapping limitations cannot substitute for a final-root comparison.

At termination, a further eight-source node maps 137 source actions to 109 actions: 87 are equal after literal renaming and 50 are rewritten. Its independent UP 1.3.0 parse fails under the default global name-uniqueness policy because type and fluent names overlap (\texttt{entity} and \texttt{location}). The same unchanged PDDL parses when that parser policy permits cross-category name reuse. This is therefore a parser interoperability diagnostic, not evidence that the PDDL is inherently invalid. The two available source traces still lack verified parameter transport. The earlier three-node frontier and this later node are separate snapshots; their counts must not be pooled.

\paragraph{Finite-state preservation and checker sensitivity.}
An authored two-source fixture covers argument reversal, subtype membership, existential and universal conditions, conditional effects, universal additions and deletions, and a foreign object. Exhaustive enumeration of its 256 source valuations checks 4,352 action applications, 2,048 applicable transitions, and 256 goal evaluations with no mismatch. Four deliberately corrupted outputs separately remove a universal-effect guard, reverse an effect polarity, reverse action arguments, and admit a foreign object into source membership. The same checker detects all four. This finite fixture tests the implementation and the checker's sensitivity to these errors; it is not a proof for arbitrary domains.

An additional finite fixture checks 14 condition forms, including negated quantifiers, implication, disjunction, equality/inequality, alternating quantifiers, variable shadowing, and quantification over an empty declared type. Across 128 valuations and 28 ground action instances, all 3,584 applicability comparisons, 1,984 applicable transitions, and 128 goal comparisons agree. Removing quantified condition membership guards while retaining parameter guards yields 1,024 mismatches. These bounded authored controls exercise formula translation and checker sensitivity; they do not extend the semantic-interface or empirical source-domain claims.

\paragraph{Variable-capture regression.}
An additional authored fixture exposed a defect in the earlier compiler: an inner universal variable shadowed an action parameter after a conditional guard was moved inward. The original one-step plan was valid but its compiled counterpart missed a goal. Fresh bound-variable renaming repairs this case. Exhaustive comparison over the fixture's 16 valuations and 32 action bindings finds eight source-atom mismatches before the repair and none afterwards. The original 256-state control and all four mutation controls also pass after the change. These controls establish their finite coverage, not a general correctness proof.

\paragraph{Semantic-negative control.}
An authored capture domain distinguishes being captured alive from being dead in a burial domain. With identical supplied initial facts and goal, the distinct-predicate baseline is unsolvable, whereas an intentionally false correspondence between captured and dead permits a two-step capture--bury plan. The plan validates in the resulting model, and both original one-step source plans still validate after translation. Thus execution preservation and plan validation can both pass despite a semantically incompatible interface. The incompatibility is known by construction, not an independently annotated real-domain label or an estimate of real-world error prevalence.

\begin{figure}[t]
\centering
\fbox{\begin{minipage}{0.92\linewidth}
\textbf{Authored source rules}\par\smallskip
\(\mathit{capture\_alive}(x):\quad \mathit{alive}(x)\longrightarrow\mathit{captured}(x)\)\\
\(\mathit{bury}(x):\quad \mathit{dead}(x)\longrightarrow\mathit{buried}(x)\)
\par\medskip
\textbf{Same supplied world and goal}\quad
\(I=\{\mathit{alive}(h)\},\quad G=\mathit{buried}(h)\)
\par\medskip
\begin{tabular}{@{}p{0.42\linewidth}p{0.50\linewidth}@{}}
Distinct predicates & False interface: \(\mathit{captured}\equiv\mathit{dead}\)\\[3pt]
No action establishes \(\mathit{dead}(h)\). & \(\mathit{capture\_alive}(h)\rightarrow\mathit{bury}(h)\)\\[3pt]
Goal is unreachable. & Two-step plan independently validates.\\
\end{tabular}
\par\medskip
\textbf{Both isolated source plans still validate after translation.}
\end{minipage}}
\caption{Execution checks cannot certify interface meaning. The false predicate correspondence creates a causal connection absent from the authored source models. Source-specific regression worlds remain separate from the composed world shown here. Static type-membership facts are omitted for readability. This is a controlled counterexample, not an estimate of error prevalence.}
\label{fig:semantic-negative}
\end{figure}

\section{Additional registered development conditions: 23 October}
These conditions extend the development evidence above; they do not replace historical runs or add independent semantic labels. Tables were generated from a hash-verified inventory of 46 reports. Repeated conditions on the same worlds are correlated and are not pooled into a larger independent benchmark.

\paragraph{Matched original worlds.}
All five runnable representations use the same 15 explicit original initial worlds and positive/negative goals. The fixed machine interfaces predate the live proposals below. Component selection retains complete source action sets and the exact problem bytes, rather than reconstructing a model from the returned plan. All 12 returned plans validate; the eight compiled plans also replay in their unchanged native models. Source-local reduction recovers the native solve coverage under this budget, while this block shows no additional original-goal coverage from interfaces.
\begin{table*}[t]
\centering
\small
\begin{tabular}{lrrrr}
\toprule
Representation & Solved & Unsolv. & Wall & Other \\
\midrule
Native source & 4 & 8 & 3 & 0 \\
Zero merge, all actions & 0 & 8 & 7 & 0 \\
Zero merge, components & 4 & 8 & 3 & 0 \\
Fixed interfaces, all actions & 0 & 8 & 7 & 0 \\
Fixed interfaces, components & 4 & 8 & 3 & 0 \\
\bottomrule
\end{tabular}
\caption{All 75 fresh searches over the same 15 known original worlds, Fast Downward 26.6 with LAMA-first, a 20-second wall/CPU limit and 4096 MiB memory limit, with one planner worker. The 15 reference-unavailable slots share one 10,800-second model-construction timeout and are not counted as planning failures. Component reduction matches native solve coverage; this block shows no source-goal benefit from the fixed unreviewed interfaces. Concurrent workstation workloads preclude statistical timing claims.}
\label{tab:matched-source-worlds-oct23}
\end{table*}


\paragraph{Model-derived mixed-source probes.}
A protocol fixed 18 directed pairs, candidate ordering, budgets and removal conditions before search. Six worlds could be generated from supported action-effect/precondition bindings. Their explicit worlds are model-derived development probes. Three support both-source necessity under whole-source removals, two do not, and one remains limited. Unsupported slots remain in the registry.
\begin{table*}[t]
\centering
\small
\begin{tabular}{llccccc}
\toprule
Producer & Consumer & Aligned & Zero & $-P$ & $-C$ & Necessary \\
\midrule
HANS\_NP & CHIC\_NP & T & T & U & T & Unresolved \\
BIRT & PUSS & S & U & U & U & Both \\
BIRT & PIGS\_NP & S & S & S & U & Not both \\
JACK\_NP & BIRT & S & U & U & U & Both \\
JACK\_NP & PIGS\_NP & S & S & S & U & Not both \\
PIGS\_NP & BIRT & S & U & U & U & Both \\
\bottomrule
\end{tabular}
\caption{All six generated worlds from 18 fixed directed source-pair slots; 12 pairs could not be constructed, retaining 48 unsupported condition slots. S: validated plan; U: solver-reported unsolvable; T: wall limit. $-P$/$-C$ remove the entire producer/consumer action set. Only three generated worlds establish both-source necessity; two remain solvable after a source removal, and one aligned case is limited. These model-derived worlds provide no independent semantic labels.}
\label{tab:composition-probes-oct23}
\end{table*}


\paragraph{Fresh output-contract pair and proposal bridge.}
A separate sequential pair attempts all 100 candidates under each schema. A common checker accepts 51 coupled outputs and zero free-field outputs, although 95 outputs on each side satisfy their respective JSON schemas. This is one paired development condition, not a semantic or repeated-run improvement estimate. Offline export retains all 234 source action variants. Only coupled output with explicit experimental role hypotheses exports five interface groups; three other export conditions yield the identical zero-merge PDDL. The resulting two unique artifacts are not four independent models.
\begin{table*}[t]
\centering
\small
\begin{tabular}{lrrrr}
\toprule
Schema policy & Checked & Invalid & Failed & Unrequested \\
\midrule
free\_fields & 0 & 95 & 5 & 0 \\
coupled\_fields & 51 & 44 & 5 & 0 \\
\bottomrule
\end{tabular}
\caption{Each condition attempted the same 100 known candidates under a 600-second request-admission/socket budget. Both produced 95 responses compliant with their respective JSON schemas, while the common contract checker accepted different counts. Checker acceptance and schema compliance are not semantic accuracy. One sequential paired run does not establish a stochastic or timing improvement.}
\label{tab:live-contract-oct23}
\end{table*}


\paragraph{New-model source replay and admission regression.}
One fixed model produced by that bridge admits six available historical native source plans, with all remaining nine source rows retained as lacking native-plan evidence. Exact full-argument inverse transport and source-state projections agree along all six traces. This adds no planner search success or interface-semantic judgment. Separately, an invalid-world control exposed a missing constant-membership admission check: a constant-bearing excluded action could affect a retained predicate. The selector now rejects such incomplete worlds instead of inserting facts. All 30 legitimate component inputs from the matched block remain accepted with byte-identical selected PDDL.
\begin{table*}[t]
\centering
\small
\begin{tabular}{lrrrr}
\toprule
Source & Steps & States & Atom checks & Mismatch \\
\midrule
ACCO & 6 & 7 & 1533 & 0 \\
BIRT & 12 & 13 & 9529 & 0 \\
CHIC\_NP & 40 & 41 & 11480 & 0 \\
HANS\_NP & 8 & 9 & 4689 & 0 \\
MUSI & 14 & 15 & 10185 & 0 \\
SOUP & 2 & 3 & 915 & 0 \\
\bottomrule
\end{tabular}
\caption{All six available historical native 60-second original-goal witnesses in one fixed live-proposed parameterized model. Nine sources without a native witness remain in the 15-source registry. No new planner calls were made. The 38,331 truth comparisons and exact full-argument inverse checks concern these correlated source-only traces, not interface meaning or mixed-source behavior.}
\label{tab:live-model-replay-oct23}
\end{table*}


\paragraph{Live-model mixed-source follow-up.}
Applying the same frozen generation rule to the new live model registers 24 directed pairs and 96 condition slots. Of 5,033 candidate records, 4,886 are excluded and 147 supported; six worlds are generated and 18 pairs remain unsupported. All six generated request objects, initial facts and goals are exactly equal to their counterparts in the older-interface cohort above. They therefore add no independent worlds. The first attempt retains 24 WSL access-denied engine errors before search. One separately labeled infrastructure retry, with identical generation bytes, model PDDL and condition registry, returns nine valid plans, 12 unsolvable reports and three wall limits. Source-removal audits again identify three both-source-necessary worlds, two still solvable after a removal and one limited aligned case. This connects the current live-proposed model to the existing mechanism probes; it establishes neither independent semantic correctness nor new benchmark coverage.

\paragraph{Fresh-reference boundary.}
The newly registered 15-source, empty-cache, 10,800-second local reference run terminates at its construction time limit with 13 completed internal nodes and no final root. Its 15 unavailable comparison slots share that single construction failure. An explicitly separate exact-request-cache condition terminates after 6,653 seconds with 14 saved internal nodes, 4,105 new model requests and a final validation error. All recorded input, implementation and cache integrity checks pass. The root candidate contains 164 actions, but strict UP parsing rejects all 15 exported problems for the same cross-category use of the name \texttt{location}. These are correlated reports of one shared interoperability failure, not 15 independent planning failures. The conditional additional-cache and root-planning gates therefore remain closed. Neither condition supplies fresh-budget reference success.

An isolated parser-policy diagnostic parses the unchanged candidate domain and all 15 problems when cross-category names are permitted. A separately registered, offline category-qualified export subsequently passes strict parsing for the same cases and inverts to the exact original syntax trees. It changes identifiers only, retains full parameter order and leaves the failed original run unchanged; no model or planner requests are made. This interoperability derivative does not establish a completed reference comparison or source preservation. Mapping inspection classifies 122 of the 234 source actions as literal-renaming matches and 112 as rewritten. All six available historical source plans remain outside the verified parameter-transport contract, with nine sources lacking a native witness; unverified transport is not an invalid-plan verdict. Intermediate-node trace counterexamples remain separate evidence.

\paragraph{Bounded posthoc transport diagnostic.}
Of the 234 root aliases, 130 retain their ordered formal signatures, 65 change arity, and 39 change formal names or order. A separately frozen structural check enumerates 176 type-compatible bijections for the latter 39: four have a unique structural match, one is ambiguous, and 34 have none. This check retains every argument and does not supply independent semantic labels. Under only unchanged signatures or these four unique maps, a further offline replay registers all 234 historical action-invocation tasks against the category-qualified derivative. Of 146 available audited source prefixes, 75 lack verified transport and 71 are admitted: 61 pass, three encounter a target-inapplicable action while the source action remains applicable, six cannot form the full source-atom projection because of target predicate types, and one raises an invalid-grounding error. Across completed comparisons, 149,117 atom truth checks have no mismatches; applicability failures and unverified errors remain in the registry. Original goals are not evaluated. These correlated diagnostic traces neither complete the failed reference condition nor establish general preservation.

An auxiliary, separately recorded public-applicability check resolves the PIGS invalid-grounding row: the first three steps execute on both sides, but the fourth source action is applicable and its mapped target action is not. Calling \texttt{apply} directly at that point raises the recorded exception. The original 234-row registry remains unchanged; this follow-up is not an additional independent task.

\paragraph{Compiler replay and admission.}
An isolated replay loads five archived compiler modules for each of the four proposal-export conditions and reproduces both PDDL and metadata bytes in all four; these conditions still contain only two unique PDDL artifacts. A separate admission check found that domain-only compiler calls could omit independent domain parsing when no problem files were supplied. Both compiler paths now require that parse explicitly. An injected checker-failure control verifies rejection, and rebuilding the same four valid export conditions retains their PDDL bytes; metadata differs only in the recorded implementation digest. These are engineering and reproducibility controls, without new inference, search, independent annotations, or benchmark instances.

Additional authored fragment controls reuse one fixed compiled model across 32 fixed-seed valuations for each of two small source models, including the all-false and all-true corners. They cover quantified conditions, disjunction, implication, equality, conditional universal and delete effects, variable shadowing, and a non-involutive ternary permutation. A first checker aborts all 64 cases because direct \texttt{apply} raises on statically false groundings; those errors are retained. A separately frozen public-applicability correction preserves all 132 fixture/model files and completes 896 applicability checks and 9,074 truth comparisons without differences, retaining 1,152 actual arguments. These are sampled engineering controls, not exhaustive coverage of the 8,192 possible valuations per source or independent semantic evidence.

\section{Additional execution diagnostics: 24 October}
\paragraph{Motivation and conditions.}
Grounding limits in existing mixed-source probes motivated two separately registered implementation conditions. Each compares 24 fixed PDDL representations of the same six known development worlds, yielding 48 strategy slots with alternating strategy order, one search worker, and a 20-second limit per process. Representation variants do not add worlds. A local reference-construction process runs concurrently; consequently, observed times do not identify causal speedups.

\paragraph{Positive-tuple joins.}
Joining positive precondition tuples retains complete actual bindings, with other conditions checked by the original simulator. The default and join strategies each return nine solved, two unsolvable and thirteen time-limited rows; all 24 paired statuses agree. Independent auditing revalidates all eighteen saved plans. The sparse authored control's smaller candidate set therefore does not translate into observed solve-coverage gain on this matrix.

\paragraph{State-local witnesses.}
The subsequent condition keeps one applicable complete binding per action and effect-parameter assignment in each state. Parameters referenced in conditional guards and effect values remain distinguished; simulated effects and quality metrics conservatively retain all parameters. The default again yields nine solved, two unsolvable and thirteen time limits. Witness selection yields ten solved, two unsolvable, eight time limits and four 5,000-state limits. All nineteen saved plans independently revalidate. The additional solved slot is the full/aligned BIRT--PIGS representation; both strategies still solve some representation of the same five of six worlds. HANS remains resource-limited. These observations provide neither new independent semantic cases nor a general performance result.
\begin{table*}[t]
\centering
\small
\begin{tabular}{lrrrrrr}
\toprule
Strategy & S & U & Time & State & Other & Worlds \\
\midrule
Default (correlated pair) & 9 & 2 & 13 & 0 & 0 & 5 \\
Positive-tuple join & 9 & 2 & 13 & 0 & 0 & 5 \\
Default (witness pair) & 9 & 2 & 13 & 0 & 0 & 5 \\
Precondition witness & 10 & 2 & 8 & 4 & 0 & 5 \\
\bottomrule
\end{tabular}
\caption{Two separately registered paired conditions, each with 48 search slots over 24 fixed PDDL representations of the same six known worlds. Every process has a 20-second limit; state limits are 5,000. S: solved, U: solver-reported unsolvable. Worlds counts a solution in any representation. All 18 and 19 saved plan records in the respective conditions independently revalidate. Defaults and worlds repeat across conditions; neither the extra witness representation nor repeated plans add independent-world coverage. Concurrent inference precludes causal timing claims.}
\label{tab:execution-diagnostics-oct24}
\end{table*}


\paragraph{Fresh reference and unavailable comparisons.}
A separately registered 18,000-second fresh reference build terminates after 18,000.031 seconds with thirteen internal nodes and no completed root. It starts with a new per-run cache and selects reversible category-qualified final export. The two-source pilot completes and strictly parses its domain and two explicit worlds, but the full run does not reach final export. Thus neither a completed fifteen-source reference nor the effect of final namespace qualification is established. Both downstream registries preserve the shared construction failure: all ninety original-world slots and all 234 given-prefix slots are unavailable, with zero planner searches and zero prefix replays. These are two linked procedures following one failure, not 324 planning failures. Historical prefix availability remains 146 audited traces and 88 missing traces. The larger construction budget and concurrent workstation loads prevent a causal comparison with the earlier 10,800-second run.
\begin{table*}[t]
\centering
\small
\begin{tabular}{>{\raggedright\arraybackslash}p{3.1cm}>{\raggedright\arraybackslash}p{2.0cm}>{\raggedright\arraybackslash}p{3.0cm}>{\raggedright\arraybackslash}p{3.0cm}}
\toprule
Stage & Units & Observed & Outcome \\
\midrule
Fresh pilot & 2 sources & 117 cache records & Complete; strict domain + 2 worlds \\
Fresh full build & 15 sources & 11166 cache records & 18,000-second time limit; no root \\
Original-world comparison & 90 slots & 0 searches & All unavailable \\
Given-prefix replay & 234 slots & 0 replays & All unavailable \\
\bottomrule
\end{tabular}
\caption{Distinct stages with different units. The pilot takes 260.579 seconds; the full build takes 18000.031 seconds and retains 13 internal nodes. Both downstream registries link to one construction timeout; their 324 unavailable slots are neither planning failures nor 324 independent tasks. Cache records count saved JSON-decoded responses, not all HTTP attempts or tokens. No cache continuation is run, and the larger budget prevents a causal comparison with the earlier 10,800-second run.}
\label{tab:reference-availability-oct24}
\end{table*}


\paragraph{Recorded workload and intermediate diagnostics.}
An offline cache audit checks all 11,166 saved request keys: 8,414 predicate checks, 2,635 action checks and 117 action updates. The adapter increments its network counter only after successful JSON decoding. Failed HTTP or JSON attempts and token use are not exhaustively logged, so the recorded responses are not an exact total inference-cost ledger. No new model or planner calls are made by the audit.

A separately bounded, posthoc audit inspects every completed internal node and reuses historical sixty-second native-plan witnesses. Across forty-four node/source occurrences, twenty-six lack native-plan evidence and thirteen have unverified parameter transport. Five mapped-plan validations yield three VALID and two INVALID records; both invalid records repeat the same BIRT native plan, while the five records use only three distinct native plans. Twelve raw internal-node domains parse and one retains a parsing error. This raw-node diagnostic is separate from the timed construction and does not establish its timeout cause. All thirteen internal PDDL hashes also occur in historical runs. Consequently, these records are correlated reproductions, not thirteen new independent models, a completed reference, or evidence that no alternative plans exist.
\begin{table*}[t]
\centering
\small
\begin{tabular}{rrlrrrr}
\toprule
Node & Sources & Raw parse & V & I & Unverified & No trace \\
\midrule
15 & 2 & VALID & 1 & 1 & 0 & 0 \\
16 & 2 & VALID & 1 & 0 & 0 & 1 \\
17 & 2 & VALID & 0 & 0 & 1 & 1 \\
18 & 2 & VALID & 0 & 0 & 1 & 1 \\
19 & 2 & VALID & 0 & 0 & 0 & 2 \\
20 & 2 & VALID & 0 & 0 & 0 & 2 \\
21 & 2 & VALID & 0 & 0 & 1 & 1 \\
22 & 3 & VALID & 0 & 1 & 1 & 1 \\
23 & 4 & VALID & 1 & 0 & 1 & 2 \\
24 & 4 & VALID & 0 & 0 & 1 & 3 \\
25 & 4 & VALID & 0 & 0 & 1 & 3 \\
26 & 7 & VALID & 0 & 0 & 4 & 3 \\
27 & 8 & ERROR & 0 & 0 & 2 & 6 \\
\bottomrule
\end{tabular}
\caption{All 13 completed internal nodes of the time-limited fresh run, not a completed root. V/I count VALID/INVALID transported historical native-plan instances under unchanged ordered formal signatures. The five validations use only three distinct native plans; both invalid records repeat BIRT. Missing mappings or traces remain unverified. All 13 node PDDL hashes also occur in historical runs, so these are correlated reproductions. The raw-node parse error is separate from the construction time limit and is not evidence that final category-qualified export failed.}
\label{tab:intermediate-nodes-oct24}
\end{table*}


\paragraph{Evidence inventory.}
The fixed reports and audits are indexed in \path{result_sources_2026_10_24_reference_v3.json}. Its 115 heterogeneous records include historical results, derived summaries and successive engineering receipts, rather than 115 independent experiments. The bounded work window and actual elapsed time are operational metadata, not evidence of scientific readiness. Independent semantic labels remain absent, and neither workload accounting nor plan validation closes that gap.

\section{Evaluation required for a submission}
\paragraph{Semantic evidence.}
Freeze candidate generation, source/context hashes, decision schemas, and evaluation splits before an independent annotation round. Annotators should see declarations and action usage but not solver outcomes. Retain equivalent, permuted-equivalent, abstraction, complementary, incompatible, and uncertain judgments separately. Report annotation coverage and disagreement; do not use model acceptance or solver success as gold labels. Current independent semantic-annotation coverage is unestablished.

\paragraph{Prepared annotation sample.}
A frozen development packet samples 100 pairs from 11,410 equal-arity cross-source pairs, within 25,868 total cross-source pairs. Sampling is deterministic within predicate-name and type-spelling strata; it includes permuted type-name signatures and unequal type names, without asserting their compatibility. Complete relevant actions retain conditional and quantified contexts. The 14,458 different-arity pairs are outside this packet's scope. No model predictions or solver outcomes appear in the packet, and no independent labels have been completed. Stratified sample proportions must not be reported as unweighted population prevalence.

The prepared agreement tool retains all six labels, reports paired coverage separately from missing judgments, and computes descriptive raw agreement and Cohen's kappa before adjudication. Argument-map agreement has its own denominator. Identical constant marginals yield undefined kappa, not an artificial perfect score. The current readiness report contains zero annotation files and null agreement scores; declared annotator independence and agreement would not themselves prove semantic truth.

\paragraph{Proposal output-contract diagnostics.}
A local Gemma 3 4B run processed all 100 packet pairs with fixed options. None passed the current structural response contract: 95 responses failed field or evidence checks, and five reached the response-token limit. Overlapping diagnostics include 75 non-equivalence labels with a non-null argument map, 42 missing two-sided evidence records, 12 non-bijective maps, and eight label/map disagreements. Offline replay reproduces all 100 result statuses after an explicit UTF-8 decoding fix. These are output-contract failures, not a measured zero semantic accuracy. A schema intervention would be a new development condition and could not by itself demonstrate improved semantic judgment.

An offline intervention audit validates 100 coupled schemas with 1,340 synthetic label/map controls. All 95 saved complete responses satisfy the original free-field schema but violate the coupled schema. This identifies a format-constraint gap; it does not measure how the model will respond under the new schema. A prepared paired-run checker requires identical model digests, prompts, options, packet contents and requests except for the schema, retaining all failures in the denominator.

A completed coupled-schema run on the same 100-pair packet yields 52 structurally checked responses, 46 responses missing evidence from both sides, and two token-truncated responses. Offline replay reproduces these statuses. After restoring the local service, a fresh free-field control on the identical packet yields zero checked, 96 invalid, and four token-truncated responses; offline replay reproduces these counts. The paired requests match except for the output schema, and both runs record the same service version and client binary hash. Thus the observed checked-output difference is 52 of 100 under this single paired development condition. It is not a repeated-run estimate, a controlled timing comparison, or a semantic accuracy result. Independent semantic correctness remains unmeasured.

\paragraph{Joint proposal consistency.}
A separate prepared check composes argument permutations across connected equivalence proposals. It withholds an entire component if a cycle assigns conflicting argument positions, two different predicates from the same source become identified, or a definite non-equivalence judgment occurs inside the component. No edge is silently discarded to repair a conflict. Exhaustive three-argument triangle controls cover all 216 combinations of edge permutations. These checks establish structural consistency only: they infer neither type correspondences nor semantic truth, and they do not approve executable alignments. The original prediction run supplies no structurally checked edges, so its empty graph is an absence of usable proposals, not successful consistency coverage. The coupled-schema run supplies seven equivalence edges in five structurally consistent components. All five components are acyclic; thus these predictions do not exercise cycle-conflict rejection. Their semantic equivalence is unverified, and none is automatically adopted.

\paragraph{Controlled comparisons.}
Compare source-isolated models, the namespaced union, a fresh UniDomain-style reference, explicit alignment, and domain-selection ablations on identical supplied worlds and goals. Separately ablate argument permutations, source-universe guards, explicit entity membership, and action-variant preservation. Report the full denominator, invalid plans, unsupported cases, and each cutoff category. Fresh reference outputs must not be replaced with a historically invalid artifact.

\paragraph{Narrative evidence and scaling.}
Add a declared candidate/ranking policy and independent assessments of event progression, entity continuity, redundancy, and coherence. Cross-source counts or source-removal necessity are diagnostics, not definitions of narrative quality. Measure compilation size, grounding overhead, planning time, and source preservation separately: membership encoding can enlarge the ground fluent space even while protecting semantics. Planner scaling analysis and broader independently sourced models remain priorities.

\section{Claim--evidence map and reviewer checklist}
\begin{description}
\item[Supported: pre-search artifact.] The entry point saves the parameterized model and manifest before subprocess search; persistent mode checks use identical domain hashes across multiple problems.
\item[Conditional: first-repeat preservation.] Full retains 499/499 available source-plan/domain transports, using 33 distinct Native plans; 401/900 slots remain unavailable. In the 19 action-merging groups, 35/35 available transports pass, but only 11/38 source-branch occurrences are exercised.
\item[Conditional: cross-source composition.] Full confirms 50/70 cross-source plans with ten origin-unknown; fixed world representatives yield 27/37 with five unknown. Predicate-only confirms 56/70 and 30/37. All comparable verdicts agree. These conditional generator counts are not a 300-group success rate.
\item[Authored four-source participation.] A separate 11-step plan uses four sources with one cross-source link and independent goal branches. Physical search is followed by original Glaive validation after an explicit apology-intention revision. This does not establish a four-source causal chain or independent story quality.
\item[Historical: execution preservation.] Complete namespaced schema recovery covers the earlier 15 sources; plan translation covers six available Fast Downward source plans. The zero-merge type-encoding replay validates all six.
\item[Fixture controls.] Synthetic tests show the intended compilation behavior and explicit rejection conditions.
\item[Semantic and narrative evidence.] No current result establishes accurate automatic real-domain alignment, superior stories, or independent-domain transfer.
\item[Contribution review: unresolved.] Establish what new scientific result goes beyond standard compilation, domain matching, and validation practice. Do not present routine encoding as novel by itself.
\item[Clarity review: addressed in this draft.] Separate source-local preservation, interface assumptions, and narrative interpretation; maintain these terms across method and experiments.
\item[Experimental strength: unresolved.] Current evidence is developmental and coverage-limited. A completed matched reference comparison and independent interface/entity evaluation remain absent.
\item[Evaluation completeness: unresolved.] Remaining repetitions, independent interface/entity annotations, broader source-trace coverage, three/four-source chain tasks, a completed matched reference, and narrative judgments remain required for their respective claims.
\item[Design soundness: partly addressed.] Type/identity counterexamples and source replays support the representation contract; mixed-source semantics and scaling still require evaluation.
\end{description}

\bibliography{references}
\end{document}
